\begin{abstract}{eBPF；开源 AI 推理框架；跨语言程序分析；调用链追踪；Ollama；HybriChain}
随着开源大语言模型在企业私有化部署中的广泛普及，以 Ollama 为代表的推理引擎采用 Go-CGO-C++ 三层混合编程与双进程解耦架构，在系统层面引入了前所未有的观测盲区：当恶意攻击者通过 API 发起目录遍历攻击或畸形载荷注入时，安全运维人员既无法追溯外部数据在推理引擎内部的完整流转路径，也无法回答该行为在内核层面触发了哪些系统调用。

本文设计并实现了 HybriChain——一个面向多语言推理引擎的跨语言调用链追踪与安全分析系统。针对跨语言符号映射难题，提出了五层符号分类体系（f$_0$–f$_4$），将来自不同命名空间（Go 运行时、CGO ABI、底层计算库）的函数符号统一映射到同一语义空间；针对追踪开销难题，设计了基于 eBPF 的探针下沉策略，仅对推理计算核心函数挂载 uprobe，将运行时开销降至最低；针对噪声数据难题，设计了五步漏斗式降噪流水线，将原始事件流压缩 2–3 个数量级；针对调用链缝合难题，提出了基于时间戳对齐的置信度缝合算法，为每条跨语言边赋予 Confirmed/Inferred/Unknown 状态；在此基础上，构建了异常行为分析引擎，支持对 futex 阻塞、异常 mmap 分配、敏感路径访问等高危行为的自动化检测。

本文以 Ollama 为典型靶点验证了 HybriChain 的有效性。实验结果表明，该系统能够从原始内核事件流中还原出完整的跨语言调用链，准确识别推理引擎中的异常行为，为开源 AI 推理软件的可观测性与安全审计提供了新的技术手段。
\end{abstract}